\section{Introduction}
\label{sec:introduction}

Many robot-learning evaluations divide interaction into short tasks and reset
the simulator after each one. This convention is useful for optimization and
comparison, but it can erase internal physical state---such as temperature or
wear---that would survive an ordinary task boundary. Removing embodiment resets
entirely creates its own exploration and learning problem
\citep{farrahi2025continuing}. We study a complementary selective-reset
setting: the arm, object, and goal reset, while a hidden actuator-health state
persists and changes the dynamics faced by later tasks.

Thermal-aware robot learning is not new. Recent quadruped work incorporates
motor-temperature models into reinforcement-learning policies and validates
thermal management on hardware \citep{qian2026thermal,wan2026residual}.
Damage-aware manipulation benchmarks also accumulate thermal and mechanical
health signals across interaction \citep{balaji2026oopsieverse}. Our question
is narrower. We deliberately use a phenomenological thermal law to isolate
what happens when health is hidden, action-driven, persistent across task
resets, and causally changes subsequent actuator dynamics. This controlled
setting is intended to expose a deployment issue, not to predict the
temperature of a physical motor.

The deployed controller has two layers. A fixed low-level policy performs the
manipulation task, while a supervisor chooses between high- and low-power modes
using a belief about the hidden thermal state. The supervisor can enlarge its
belief by an uncertainty margin before making that choice. Such modular
uncertainty-aware filtering has strong precedent, including belief-space and
latent safety filters with substantially stronger guarantees than those studied
here \citep{vahs2024risk,seo2025latent,kwon2026adaptive}. We therefore do not
claim a new safety-filter architecture. Instead, we ask whether a fixed margin
improves lifetime utility under deployment shifts, whether that numerical
margin transfers between robot-control tasks, and what a target-specific
calibration stage actually recovers.

We answer these questions in Pusher and Reacher under matched selective-reset
semantics and a shared thermal law. Each independent evaluation unit is a
20-task lifetime. Sensor-noise, cooling, physical-shock, and combined shifts
are evaluated using paired fresh lifetimes. The experimental sequence keeps
the inherited cross-task test separate from a later, development-only
calibration and its disjoint fresh test set.

The results provide three main findings. First, a fresh Pusher factorial
experiment attributes the robust benefit primarily to the uncertainty margin:
without the learned residual, the margin improves target-shift mean utility by
$+0.965$ reward per task (95\% bootstrap CI $[0.733,1.196]$). Second, the same
margin retains a positive mean effect on Reacher ($+0.721$,
$[0.474,0.978]$) but reaches a 2.3\% maximum trip rate and therefore fails the
prespecified 2.0\% tolerance. Third, a separately frozen calibration using only
Reacher development lifetimes selects a more conservative margin; on 100 new
lifetimes it reaches 1.6\% while preserving a positive mean effect ($+0.692$,
$[0.434,0.962]$). Its mean utility is not detectably different from the
inherited setting. Calibration is therefore credited with tolerance recovery,
not reward superiority.

These are mean risk-sensitive utility results, not universal dominance. Only
52 of 100 calibrated Reacher lifetime effects are positive, and the exact sign
test is not significant. Pusher reward decomposition further shows that fewer
thermal trips compensate for lower immediate task return, so the conclusion
depends on trips carrying moderate or high cost. Together, the evidence
supports a practical and testable workflow---calibrate the uncertainty margin
on target-task development lifetimes, freeze it, and evaluate it once---while
leaving hardware validity, formal safety, and general algorithmic superiority
open.
